\begin{exercise}[Détermination expérimentale du c{\oe}fficient de dilatation thermique]
La structure de l'alliage, de composition définie \ce{FeSi}, a été étudiée par diffraction  de neutrons sur poudre entre 4 et \SI{1173}{\kelvin}. Aucune transition polymorphique n'a été observée, la structure restant cubique de groupe d'espace $P2_13$ sur toute la plage de température.

\begin{questions}
\question Selon vous, quel avantage comporte l'utilisation de la diffraction de neutrons par rapport à la diffraction d'électrons dans ce système ?

\`{A} pression constante, le paramètre de maille a été relevé à différentes températures, permettant le calcul du volume de maille.

\begin{center}
\begin{minipage}{0.5\textwidth}
\includegraphics[width=\textwidth]{/home/jeff/Documents/IUTOrsay/1-Cours/4p13-Alliages/Cours-TD/Exos/CP/Coef-ThermElas-Ex01/figures/Fig20}
\end{minipage}\begin{minipage}{0.5\textwidth}
\captionof{figure}{\'{E}volution du volume de la maille cubique en fonction de la température en échelle lin-log. Le point $(300 ~ \si{\kelvin}, 90.22 ~ \si{\angstrom\cubed})$ a servi de point de référence à la courbe.}
Le modèle a pour équation : 
\begin{align*}
\frac{V}{V_r} = \exp \left( 4.85 \cdot 10^{-5} \cdot (T-300) \right)
\end{align*}
\end{minipage}
\end{center}

Sachant que le système a une température de Debye de \SI{445}{\kelvin}, on va chercher à évaluer l'évolution du volume pour $T > T_r = 300 ~ \si{\kelvin}$, sachant que ce cristal a un volume de $V_r = V(\SI{300}{\kelvin}) = 90.22(1) \si{\angstrom\cubed}$.

\question \'{E}crire le développement limité au premier ordre du volume en fonction de la température.

\question Faire apparaître le c{\oe}fficient de dilatation thermique $\alpha$.

On va supposer que le solide \ce{FeSi} est faiblement dilatable sur la plage de température d'étude, ce qui permet d'écrire que $\alpha = \alpha_0 = cste$.

\question Montrer que le volume du solide peut s'écrire comme : 
\begin{align*}
V = V_r \exp \left( \alpha_0 (T - T_r) \right)
\end{align*}

\question En déduire la valeur de $\alpha_0$ de \ce{FeSi} du modèle simple utilisé.

Un modèle thermodynamique plus complexe, appelé modèle de Grünseisen, a permis d'obtenir l'évolution du c{\oe}fficient de dilatation thermique sur toute la plage de température. 

\begin{center}
\begin{minipage}{0.5\textwidth}

\includegraphics[width=\textwidth]{/home/jeff/Documents/IUTOrsay/1-Cours/4p13-Alliages/Cours-TD/Exos/CP/Coef-ThermElas-Ex01/figures/Fig21}
\end{minipage}\begin{minipage}{0.5\textwidth}
\captionof{figure}{\'{E}volution du c{\oe}fficient de dilatation thermique à pression constante de \ce{FeSi} entre 4 et \SI{1173}{\kelvin}.}
\end{minipage}
\end{center}

\question Commenter l'allure de la courbe du c{\oe}fficient de dilatation thermique de l'alliage \ce{FeSi}.
\end{questions}
\end{exercise}